% Papier à en-tête ICube
% Version du 3 février 2013
% Infos et bugs : vincent.mazet@unistra.fr

\documentclass[a4paper,11pt,twosided]{letter}

% Paquets généraux
\usepackage[utf8]{inputenc}     % Codage du fichier TeX
\usepackage[T1]{fontenc}        % Codage du fichier PDF (les paquets cm-super ou lmodern doivent être installés pour obtenir un bon résultat)
\usepackage{graphicx}           % Logos et autres illustrations
\usepackage{fancyhdr}           % En-têtes et pieds de page
\usepackage{color}              % Couleur
\usepackage{ifthen}             % Condition
\usepackage{helvet}             % Police Helvetica
\usepackage[french]{babel}      % Courrier en français
% \usepackage[english]{babel}   % Courrier en anglais

% Papier à en-tête ICube
% Version du 3 février 2013
% Infos et bugs : vincent.mazet@unistra.fr

% Mise en page
\voffset -1in
\topmargin 10mm
\headheight 5mm
\headsep 35mm
\textheight 210mm
\footskip 20mm
\hoffset -1in
\oddsidemargin 50mm
\evensidemargin 50mm
\textwidth 140mm
\marginparsep -180mm
\marginparwidth 35mm
\parindent 0mm

% En tête et pieds de page
\pagestyle{fancy}
\renewcommand{\headrulewidth}{0pt}

% Police sans-serif
\renewcommand{\familydefault}{\sfdefault}

% Couleurs
\definecolor{bleu}{RGB}{0,94,168}
\definecolor{gris}{RGB}{88,88,90}

% Logos
\newcommand{\icubesign}{\includegraphics[width=40mm]{icube-sign.eps}\vspace*{-20mm}}
\newcommand{\icube}{\includegraphics[width=40mm]{icube.eps}\vspace*{-12mm}}

% Marge (expéditeur, suivi, adresse)
\newcommand{\marge}{%
  \vspace*{-225mm}\hspace*{-40mm}\vbox{
    \vtop to  45mm {\vfil\usebox{\exped}\vfil}
    \vbox to  30mm {\vfil}
    \vbox to  30mm {\vfil\usebox{\suiv}\vfil}
    \vbox to 100mm {\vfil}
    \vbox to  20mm {\vfil\usebox{\adress}}
  }
}

% Tutelles
\newcommand{\tutelles}{
  \begin{tabular}{@{}r@{}}%
  \scriptsize Sous la co-tutelle de \\%
  \hline%
  ~\\[-3mm]%
  \includegraphics[height=1cm]{cnrs.eps}\quad%
  \ifthenelse{\value{engees}=1}{\includegraphics[height=1cm]{engees.eps}\quad}{}%
  \ifthenelse{\value{insa}=1}{\includegraphics[height=1cm]{insa.eps}\quad}{}%
  \includegraphics[height=1cm]{unistra.eps}%
  \end{tabular}
}

% Éléments de la lettre
\lhead{\hspace*{-40mm}\ifthenelse{\value{page}=1}{\icubesign}{\icube}}      % Logo ICube
\chead{}                                                                    %
\lfoot{\ifthenelse{\value{page}=1}{\marge}{}}                               % Expéditeur, suivi, adresse
\cfoot{}                                                                    %
\rfoot{\ifthenelse{\value{page}=1}{\tutelles}{\raisebox{-5mm}{\thepage}}}   % Tutelles

% Expéditeur
\newsavebox{\exped}
\newcommand{\expediteur}[3]{\savebox{\exped}{\parbox{50mm}{%
  \textcolor{bleu}{#1 \textbf{#2}} \\[1ex]\scriptsize#3}}}

% Suivi de l'affaire
\newsavebox{\suiv}
\newcommand{\suivi}[1]{\savebox{\suiv}{\parbox{35mm}{\scriptsize%
  \textcolor{bleu}{Affaire suivie par~:} \\[.5ex]#1}}}

% Adresse du laboratoire
\newsavebox{\adress}
\newcommand{\adresse}[1]{\savebox{\adress}{\parbox{50mm}{%
  \scriptsize\color{gris}\textbf{Laboratoire ICube --- UMR 7357}\\#1}}}

% Autres commandes
\newcommand{\lieudate}[1]{\rhead{\ifthenelse{\value{page}=1}{#1}{}}}
\newcounter{engees}
\newcommand{\logoengees}{\setcounter{engees}{1}}
\newcommand{\nologoengees}{\setcounter{engees}{0}}
\newcounter{insa}
\newcommand{\logoinsa}{\setcounter{insa}{1}}
\newcommand{\nologoinsa}{\setcounter{insa}{0}}
\newcommand{\destinataire}[1]{\vspace*{-20mm}\vbox to 4cm {\vfil \hfill\hbox to 100mm{\parbox{90mm}{#1}}\vfil}\vspace*{20mm}}
             % Style de la lettre à en-tête

\begin{document}

% Lieu et date
\lieudate{Illkirch, le \today}

% Tutelles
\logoengees         % Logo ENGEES
% \nologoengees     % Pas de logo ENGEES
\logoinsa           % Logo INSA
% \nologoinsa         % Pas de logo INSA

% Expéditeur
\expediteur{Vincent}{Mazet}{%
  Maître de conférences, \\
  Tél~: +33 (0)3 68 85 44 91 \\
  vincent.mazet@unistra.fr
}

% Affaire suivie par...
\suivi{%
  Prénom Nom \\
  Tél~: \\
  Mél~: \\
  Réf.
}

% Adresse du laboratoire
\adresse{%
  Télécom Physique Strasbourg \\
  300 bd Sébastien Brant -- BP 10413 \\
  F--67412 Illkirch cedex \\
  http://icube.unistra.fr
}

% Destinataire
\destinataire{%
  Laboratoire ICube \\
  300 boulevard Sébastien Brant \\
  BP 10413 \\
  67412 Illkirch Cedex -- France
}

% Objet
Objet~: rédaction d'une lettre avec papier à en-tête en \LaTeX{}

\bigskip\bigskip

Madame, Monsieur,

\bigskip\bigskip

la version \LaTeX{} du papier à en-tête du laboratoire ICube est configurable avec les éléments suivants~:
\begin{itemize}
  
  \item \texttt{$\backslash{}$lieudate\{\dots\}} permet d'indiquer
  le lieu et la date de rédaction en haut à droite sur la première page~;

  \item \texttt{$\backslash{}$expediteur\{\textit{Prénom}\}\{\textit{Nom}\}\{\dots\}}
  s'écrit dans la marge et identifie l'auteur de la lettre. Indiquez vos prénom, nom et coordonnées~;

  \item \texttt{$\backslash{}$suivi\{\dots\}} permet de donner des informations sur la personne suivant le dossier~;

  \item \texttt{$\backslash{}$adresse\{\dots\}} contient l'adresse du laboratoire,
  à personnaliser suivant le site où on se trouve~;

  \item \texttt{$\backslash{}$destinataire\{\dots\}} correspond au destinataire du courrier~;

  \item ICube possède quatre tutelles. Le CNRS et l'Université de Strasbourg doivent apparaître.
  En revanche, l'ENGEES et l'INSA sont optionnel~;
  on peut les afficher avec les commandes \texttt{$\backslash$logoengees} et \texttt{$\backslash$logoinsa}
  ou les faire cacher avec les commandes \texttt{$\backslash$nologoengees} et \texttt{$\backslash$nologoinsa}.

\end{itemize}

\bigskip

Ce document nécessite les logos des tutelles (fichiers cnrs.eps, engees.eps, icube.eps, icube-sign.eps, insa.eps et unistra.eps)
et les paquets suivants~:
\begin{itemize}
  
  \item fontenc, pour créer des fichiers PDF avec des polices correctes
  (les paquets cm-super ou lmodern doivent par ailleurs être installés pour obtenir un bon résultat)~;

  \item graphicx pour afficher des images (en l'occurrence, les logos)~;

  \item fancyhdr, pour pouvoir modifier les en-têtes et pieds de page~;

  \item color, pour prendre en charge les couleurs de la charte graphique~;

  \item ifthen, pour gérer des conditions~;

  \item helvet pour écrire avec la police Helvetica~;

  \item babel, avec l'option french ou english pour prendre en comptes les particularités de la langue.

\end{itemize}
Le fichier icube-style.tex est indispensable~: il contient les commandes nécessaires à la compilation.

\bigskip

Les logos inclus sont au format EPS. De ce fait, la compilation est à faire avec les commandes \texttt{latex}, \texttt{dvitops} et \texttt{pstopdf}.

\bigskip

Ce document correspond à la version du \today.
Pour toute information supplémentaire, ou pour faire remonter une correction, veuillez m'envoyer un courriel à vincent.mazet@unistra.fr.

\bigskip

Veuillez agréer, Madame, Monsieur, mes salutations distinguées.

\bigskip\bigskip

\hfill Vincent Mazet

\end{document}
